% !TeX root = ../../Book3.tex
% 纲要标题：### 19.5*　体积约束与常平均曲率
% 纲要位置：几何学计划纲要.md，第 2896 行。

% 纲要内容开始
% 
% > 写作范围与外部输入：规划任务：以闭、双侧、嵌入并围成区域的光滑超曲面为主范围，固定环境定向、单位法向与平均曲率归一化，完整推导体积约束的一阶和二阶变分、常平均曲率条件与约束稳定性。用第16章的球面谱或相应已证不等式完成欧氏球面的保体积稳定模型。一般 CMC 存在、分类及 min-max 构造准确外引；不由无约束稳定性替代保体积稳定性。
% 
% * 围成区域、法向与体积变分
% * 体积约束与常平均曲率方程
% * 保体积变分的二阶公式
% * 球面模型与约束稳定性
% * 存在、分类与边界条件的适用范围
% 
% ---
% 
% 纲要内容结束
% 正文状态：未撰写；本文件目前是框架。
% 规划任务：以闭、双侧、嵌入并围成区域的光滑超曲面为主范围，固定环境定向、单位法向与平均曲率归一化，完整推导体积约束的一阶和二阶变分、常平均曲率条件与约束稳定性。用第16章的球面谱或相应已证不等式完成欧氏球面的保体积稳定模型。一般 CMC 存在、分类及 min-max 构造准确外引；不由无约束稳定性替代保体积稳定性。

\OptionalSection{体积约束与常平均曲率}{g3:sec:1905}

待撰写

\subsection{围成区域、法向与体积变分}
\label{g3:subsec:1905:01}
待撰写

\subsection{体积约束与常平均曲率方程}
\label{g3:subsec:1905:02}
待撰写

\subsection{保体积变分的二阶公式}
\label{g3:subsec:1905:03}
待撰写

\subsection{球面模型与约束稳定性}
\label{g3:subsec:1905:04}
待撰写

\subsection{存在、分类与边界条件的适用范围}
\label{g3:subsec:1905:05}
待撰写
